problem:  
Let \( M^n \) be a compact, simply connected Riemannian manifold with nonnegative sectional curvature and an effective, isometric torus action \( T^k \) of maximal symmetry rank \( k = \lfloor n/2 \rfloor \). Suppose the orbit space \( Q = M/T^k \) is an Alexandrov space with curvature \(\ge 0\), and denote by \( \mathcal{S} \) its orbit-type stratification determined by isotropy subgroups of \(T^k\).  

A fundamental open issue in transformation geometry is whether the metric geometry of \(Q\), combined with the isotropy weight data along its strata, suffices to determine \(M\) up to equivariant diffeomorphism. Current rigidity results (e.g., in homogeneous or toric settings) affirm this only under strong assumptions of global symmetry or convexity.  

**Problem:**  
Formulate the following two-part rigidity question precisely within the setting of stratified Alexandrov spaces:  
1. Consider the moduli space \(\mathcal{A}(Q, \mathcal{S})\) of all Alexandrov metrics on a fixed stratified space \((Q,\mathcal{S})\) arising as orbit quotients of compact nonnegatively curved Riemannian manifolds with \(T^k\)-actions of maximal rank and fixed isotropy representation data along each stratum.  
   - Is \(\mathcal{A}(Q, \mathcal{S})\) **discrete** up to metric isomorphism?  
   - Equivalently, do sufficiently small Gromov–Hausdorff perturbations of the metric on \(Q\) within the class preserving curvature \(\ge 0\) and the isotropy weight data always arise as orbit spaces of manifolds \(M'\) that are equivariantly diffeomorphic to \(M\)?  
2. If not discrete, can one produce a continuous Alexandrov deformation family \(Q_t\) (compatible with \(\mathcal{S}\) and isotropy weights) corresponding to non–equivariantly-diffeomorphic total spaces \(M_t\)? That is, can **topological variation of the total space** occur despite first-order Alexandrov rigidity of the quotient geometry?  

Why is it a "good" problem:  
- **Profound insight:** This question refines the intuitive “deformation rigidity” idea into a well-posed framework using the Gromov–Hausdorff topology and fixed orbit-type stratification. It probes whether infinitesimal Alexandrov rigidity extends to global diffeomorphism rigidity of the corresponding manifolds—an analogue of Mostow-type or shape rigidity for spaces with nonnegative curvature and maximal torus symmetry.  
- **Bridge between fields:** The problem unites **Riemannian geometry** (nonnegative curvature, torus symmetry rank rigidity), **Alexandrov geometry** (metric and curvature convergence under Gromov–Hausdorff topology), and **equivariant topology** (isotropy representations and orbit-type stratifications). Any progress requires new synthesis across these domains—linking stability results in metric geometry to diffeomorphism classification under symmetry constraints.  
- **Extreme simplicity:** Conceptually it asks a succinct, intuitive question: *does the quotient geometry of a maximally symmetric, nonnegatively curved manifold determine the underlying manifold, even up to small deformations?* This clarity conceals the full depth of the problem, as “small” metric perturbations in Alexandrov sense interface with highly nonlinear structures in the manifold topology.  
- **Open and deep:** There is no known general rigidity or deformation theory for Alexandrov orbit spaces with fixed stratification and isotropy data. Establishing discreteness or detecting continuous moduli would reveal unknown layers of metric rigidity underlying nonnegative curvature with symmetry and would pioneer a new direction in “stratified Alexandrov deformation theory.”